\section{Calendar order on the whole line}
\label{sec:calendar}

One-expiry validity removes butterfly arbitrage.  A surface must also order
expiries: moving expiry later cannot make a normalized call less valuable when
the underlying forward-normalized price is a martingale.  Because every LQD
slice defines a price at every strike, the requirement applies on the whole
strike line, not only where two expiries happen to have quotes.

\subsection{The calendar inequality}

\begin{lemma}[Equal-log-moneyness calendar order]
\label{lem:caljensen}
Assume deterministic rates and carry.  Let $F_{0,t}$ be the forward curve and
$M_t=S_t/F_{0,t}$ a positive martingale.  For $T_1<T_2$ and every
$k\in\R$,
\begin{equation}
c_1(k)\le c_2(k).
\label{eq:calreq}
\end{equation}
\end{lemma}

\begin{proof}
With $Y_i=M_{T_i}$, conditional Jensen gives
\[
\begin{aligned}
c_2(k)
&=\E[(M_{T_2}-e^k)^+]\\
&\ge\E[(\E[M_{T_2}\mid\mathcal F_{T_1}]-e^k)^+]
=\E[(M_{T_1}-e^k)^+]=c_1(k).
\end{aligned}
\]
\end{proof}

Equal $k$ means equal strike divided by forward, so the cash strikes are
$K_i=F_i e^k$ and are generally different.  The two laws of $Y_i$ have the
same mean, while the farther law is more dispersed in the precise sense of
convex order.  If carry is stochastic, the numeraire and the objects compared
must be specified before using this relation.

Black call price is strictly increasing in $w$ at every finite $k$.  Hence
\begin{equation}
c_1(k)\le c_2(k)\quad\Longleftrightarrow\quad
w_1(k)\le w_2(k).
\label{eq:calendarvarianceequiv}
\end{equation}
Thus the visual desk rule for a stacked total-variance chart is exact:
farther-expiry curves must lie above nearer-expiry curves at every
log-moneyness.

\subsection{Move the calendar test from strike to rank}

Calls from two expiries reach a given strike at different ranks.  Directly
comparing them therefore mixes the two transport maps.  The ledger places all
expiries on the same probability-rank axis.  The following conjugacy shows
that no information is lost by making this change.

\begin{lemma}[Ledger--call conjugacy]
\label{lem:conjugacy}
For every admissible slice,
\begin{equation}
\begin{aligned}
c(k)&=\max_{z\in[-\infty,\infty]}
\{G(z)-e^k(1-\Lambda(z))\},\\
G(z)&=\min_{k\in[-\infty,\infty]}
\{c(k)+e^k(1-\Lambda(z))\}.
\end{aligned}
\label{eq:conjugacy}
\end{equation}
The optimizers are $z=z_k$ and $k=x(z)$.
\end{lemma}

\begin{proof}
For fixed $k$,
\[
G(z)-e^k(1-\Lambda(z))
=\int_z^\infty(e^{x(t)}-e^k)\rho(t)\,\dd t.
\]
The integrand is negative before $z_k$ and positive after it.  Starting the
integral at $z_k$ maximizes the expression and gives the call.  The first
identity implies
$c(k)+e^k(1-\Lambda(z))\ge G(z)$, with equality at $k=x(z)$, which gives the
second identity.
\end{proof}

\begin{theorem}[Calendar order is ledger order]
\label{thm:ledgerorder}
For two admissible mean-one slices,
\begin{equation}
c_1(k)\le c_2(k)\ \text{for every }k\in\R
\quad\Longleftrightarrow\quad
G_1(z)\le G_2(z)\ \text{for every }z\in\R.
\label{eq:ledgerorder}
\end{equation}
\end{theorem}

\begin{proof}
If $G_1\le G_2$, take maxima in the first line of
\eqref{eq:conjugacy}.  If $c_1\le c_2$, take minima in the second line.
\end{proof}

The theorem says that checking all strikes is exactly the same as checking all
ranks.  The ledger axis is common across expiries and the slice builder already
stores it, making it the natural coordinate for both calibration constraints
and the final certificate.  Ordered marginals are compatible with at least
one martingale by the theorems of Strassen and Kellerer
\citep{Strassen1965,Kellerer1972}.  The marginal surface does not select that
martingale dynamics.

\subsection{Full-line constraints in calibration}
\label{sec:globalcalendarfit}

A constraint restricted to the common quoted span controls an interpolant
there; it says nothing about the published extrapolation.  That restriction is
appropriate only for a model deliberately left undefined outside the quoted
span.  An LQD slice is defined on all of $\R$, so the accepted maturity stack
imposes
\begin{equation}
G_{j+1}(z)-G_j(z)\ge0
\qquad\text{for every }z\in\R
\label{eq:globalledgerconstraint}
\end{equation}
for every adjacent pair of expiries.

The expiries are fitted jointly because changing either member of an adjacent
pair changes the calendar gap.  The data term is the sum of the slice
objectives.  The inequalities \eqref{eq:globalledgerconstraint} are hard
constraints, not residuals whose violation can be traded against a slightly
better quote fit.  Standard constrained solvers use the analytic ledger
sensitivities from \eqref{eq:pass3}; a least-squares implementation may use the
active-set exchange described below.  In either case, acceptance requires the
full-line certificate.

Global control necessarily couples unquoted wings across maturity because
those wings belong to the published surface.  Choose the tail policy for the
stack before fitting.  With common tail exponents, eventual calendar order
requires
\begin{equation}
\lambda_{-,j+1}\ge\lambda_{-,j},
\qquad
\lambda_{+,j+1}\ge\lambda_{+,j},
\label{eq:tailscalecalendar}
\end{equation}
apart from exact ties, where the next asymptotic constant decides the order.
These inequalities are imposed in the endpoint chart.  If exponents vary by
expiry, the farther law cannot have a lighter asymptotic tail: on either side
$\alpha_{\pm,j+1}>\alpha_{\pm,j}$ forces a calendar reversal sufficiently far
out.  Common $\alpha_\pm$ across a maturity stack is the simpler policy.

If the hard constraints are inconsistent with the quote bands, no law stack
can satisfy both those bands and the selected tail policy.  The meaningful
diagnostic is the smallest quote-band relaxation needed for feasibility.
Shrinking the calendar domain would merely hide the inconsistency outside the
chosen interval.

\subsection{A complete calendar certificate}
\label{sec:calendarcertificate}

The all-line constraint still appears infinite-dimensional.  The transport
geometry reduces it to finitely many evaluations.  Let
\[
H(z)=G_2(z)-G_1(z).
\]
At both endpoints $H$ tends to zero.  Its derivative has a simple sign:
\begin{equation}
H'(z)=\rho(z)\{e^{x_1(z)}-e^{x_2(z)}\}.
\label{eq:calgapderivative}
\end{equation}
Thus the ledger gap can turn only where the two quantile curves cross.  Between
two consecutive crossings it is monotone.  Every possible negative minimum
therefore occurs at a crossing or in an analytic tail.

This gives a finite certificate for the stored surface:
\begin{enumerate}
\item isolate every root of $x_{j+1}-x_j$ on the central grid; for cubic
Hermite transports this is a cubic root problem on each segment;
\item evaluate $G_{j+1}-G_j$ at those roots and at the two grid boundaries;
\item continue both transports with \eqref{eq:rightcontinuation} and its left
analogue, isolate any tail roots, and evaluate the ledger gap there;
\item check the limiting tail order from the exponents, endpoint coefficients,
and, in a tie, the asymptotic constants.
\end{enumerate}
Nonnegative values at all these candidate minima prove
\eqref{eq:globalledgerconstraint} on the whole real line.  This is a finite
certificate of a continuum property; no arbitrary strike cutoff is involved.

During calibration, the current violating minimum is added to the active
constraint set.  This exchange step is repeated until the certificate finds
no negative minimum.  In practice the active ranks are few because adjacent
smooth quantile curves cross only a few times.

\begin{figure}[htbp]
\centering
\paperfig{0.94\textwidth}{fig_calendar_global.pdf}
\caption{A globally ordered synthetic LQD maturity stack.  Left: total
variance at four expiries; the curves do not cross.  Right: adjacent ledger
gaps on the common rank axis; each remains nonnegative and returns to zero at
both endpoints.  The displayed curves illustrate the order; the quantile-root
test certifies it on the full line.}
\label{fig:calendarglobal}
\end{figure}

At extreme strikes, both Black price and vega can underflow while the implied
variance remains finite.  A publication chart therefore evaluates remote
wings from \eqref{eq:leeclosed} or \eqref{eq:rightsublinearwing}, matched to
the central inversion, rather than inverting numerically zero prices.  The
chart then displays the globally certified surface instead of numerical noise
from an ill-conditioned inverse.
